\documentclass[a4paper,11pt]{article}
\usepackage{ctex}
\usepackage{amsmath, amssymb}
\usepackage{geometry}
\usepackage{booktabs}
\geometry{left=2.5cm, right=2.5cm, top=2.5cm, bottom=2.5cm}

\begin{document}

\begin{center}
    \Large\textbf{上机作业：条件数估计与计算解精度分析}
\end{center}

\section{实验目的与算法原理}
本实验的主要任务是将矩阵 1-范数条件数估计的优化算法（算法 2.5.1）编制为通用子程序，并将其应用于两类经典线性系统的后验精度评估中。

\subsection{无穷范数条件数的估计转化}
题目要求估计矩阵 $A$ 的无穷范数条件数 $\kappa_\infty(A) = \|A\|_\infty \|A^{-1}\|_\infty$。其中 $\|A\|_\infty$ 可通过计算绝对值行和的最大值直接求得。
为了利用给定的算法 2.5.1（用于估计 1-范数），我们依据矩阵范数的对偶性质：
\begin{equation}
\|A^{-1}\|_\infty = \|(A^{-1})^T\|_1 = \|(A^T)^{-1}\|_1
\end{equation}
构造辅助矩阵 $B = (A^T)^{-1}$。在算法迭代中：
\begin{itemize}
    \item 计算 $w = Bx$ 转化为求解线性方程组 $A^T w = x$；
    \item 计算 $z = B^T v$ 转化为求解线性方程组 $A z = v$。
\end{itemize}
该转换避免了显式求逆矩阵，通过复用高斯消去法即可高效、稳定地获得条件数的近似值。

\section{计算任务与结果分析}

\subsection{任务 1：Hilbert 矩阵的条件数估计}
应用编写的子程序，估算 5 到 20 阶 Hilbert 矩阵的无穷范数条件数，部分代表性结果如表 1 所示。

\begin{table}[h]
\centering
\caption{Hilbert 矩阵无穷范数条件数估计}
\begin{tabular}{ccc}
\toprule
阶数 $n$ & 估算的 $\kappa_\infty(H_n)$ & 真实条件数参考值 \\
\midrule
5  & $9.4366 \times 10^{5}$  & $9.4366 \times 10^{5}$ \\
10 & $3.5353 \times 10^{13}$ & $3.5353 \times 10^{13}$ \\
13 & $8.5317 \times 10^{17}$ & $8.5350 \times 10^{17}$ \\
15 & $1.0915 \times 10^{18}$ & $1.0975 \times 10^{18}$ \\
20 & $2.4542 \times 10^{19}$ & $3.0396 \times 10^{19}$ \\
\bottomrule
\end{tabular}
\end{table}

\textbf{实验现象讨论}：
实验数据清晰地展示了 Hilbert 矩阵在条件数上的极端病态性。
值得注意的是，在 $n \le 12$ 时，优化算法的估计值与精确值高度一致；但当 $n = 13$ 时（此时条件数突破 $10^{16}$），两者开始出现偏差。这一现象在理论上是必然的：由于计算机双精度浮点数（Float64）的机器精度约为 $2.2 \times 10^{-16}$，当系统条件数超过 $10^{16}$ 阈值时，问题本身的数值求解已经损失了全部有效数字。这导致估计算法中求解 $A^T w = x$ 的过程本身已被巨大的截断误差掩盖，从而使得 $n \ge 13$ 时的估计值发生理论性偏移。

\subsection{任务 2：列主元高斯消去法的极端不稳定性}
针对特定结构的矩阵 $A_n$（主对角线与最后一列为 1，其余下三角区域为 -1），随机生成精确解 $x$ 构造右端项 $b$，使用列主元高斯消去法求解得到 $\hat{x}$。计算其真实相对误差与后验精度界（基于公式 $\kappa_\infty(A) \frac{\|b - A\hat{x}\|_\infty}{\|A\|_\infty \|\hat{x}\|_\infty}$），结果如表 2 所示。

\begin{table}[h]
\centering
\caption{特定矩阵 $A_n$ 的计算解精度估计对比}
\begin{tabular}{cccc}
\toprule
阶数 $n$ & 无穷范数条件数估计 & 后验精度估计界 & 真实相对误差 \\
\midrule
5  & $5.0000$ & $2.4172 \times 10^{-16}$ & $1.2086 \times 10^{-16}$ \\
10 & $1.0000 \times 10^{1}$ & $6.3333 \times 10^{-15}$ & $3.5984 \times 10^{-15}$ \\
15 & $1.5000 \times 10^{1}$ & $6.0010 \times 10^{-13}$ & $4.2443 \times 10^{-13}$ \\
20 & $2.0000 \times 10^{1}$ & $3.1066 \times 10^{-12}$ & $2.4456 \times 10^{-12}$ \\
25 & $2.5000 \times 10^{1}$ & $4.7564 \times 10^{-10}$  & $1.9048 \times 10^{-10}$  \\
30 & $3.0000 \times 10^{1}$ & $4.0112 \times 10^{-9}$  & $1.8788 \times 10^{-9}$  \\
\bottomrule
\end{tabular}
\end{table}

\textbf{实验现象讨论}：
该矩阵 $A_n$ 是著名 Wilkinson 矩阵的一个变体。由表 2 可观察到一个经典的反直觉现象：
\begin{enumerate}
    \item \textbf{方程组高度良态}：估算出的 $\kappa_\infty(A_n)$ 精确等于阶数 $n$。矩阵条件数极小，且仅呈线性缓慢增长，系统不存在病态敏感性。
    \item \textbf{误差出现指数级增长}：尽管方程组十分良态，但当 $n$ 增加到 30 时，计算解的相对误差却扩大到了 $10^{-9}$ 量级，显著丢失了有效数字。
\end{enumerate}

其根本原因在于算法自身的不稳定性。在执行列主元高斯消去法时，该矩阵天然满足列主元条件（无需任何行交换）。然而在消元过程中，矩阵最后一个元素的数值会以 $1 \to 2 \to 4 \dots \to 2^{n-1}$ 的规律发生指数级增长。当 $n=30$ 时，高达 $2^{29}$ 的增长因子（Growth Factor）使得极其微小的机器舍入误差在消元步骤中被放大了上亿倍。

此外，表 2 中后验精度估计界与真实误差量级高度吻合，且通常提供了安全可靠的上限，这验证了利用范数条件数结合残差来评估计算解精度的有效性。

\end{document}